\documentclass[10pt,a4paper]{amsart}
\usepackage[margin=19mm]{geometry}
\usepackage{amsmath,amssymb,mathtools}
\usepackage[scr=rsfs,cal=euler]{mathalfa}
\usepackage{xfrac}
\usepackage[unicode,hidelinks,bookmarks=false]{hyperref}
\newcommand{\ie}{i.e.\ }
\newcommand{\cf}{cf.\ }
\newcommand{\resp}{respectively\ }
\newcommand{\Subsection}{\subsection}
\providecommand{\nobreakdash}{\nobreak\hbox{-}}
\setlength{\emergencystretch}{2em}
\multlinegap0pt
\usepackage{url,color,xurl,tikz,color,soul}
\usepackage{amsaddr}
\usepackage{verbatim}
\newtheorem{theorem}{Theorem}[section]
\newtheorem{proposition}[theorem]{Proposition}
\newtheorem{lemma}[theorem]{Lemma}
\newtheorem{corollary}[theorem]{Corollary}
\newtheorem{conjecture}[theorem]{Conjecture}
\newtheorem{definition}{Definition}[section]
\newtheorem{example}[theorem]{Example}
\newtheorem{remark}[theorem]{Remark}
\begin{document}
\title[On smooth elements in arithmetic semigroups]{On smooth elements in arithmetic semigroups}
\author{Giovanna De Lauri, Dan Fretwell, Jack Ye}
\date{}
\maketitle
\noindent\emph{Source and licence.} On smooth elements in arithmetic semigroups. Version arXiv:2609.29748v1 (CC BY 4.0). This material is distributed under the Creative Commons Attribution 4.0 International licence, http://creativecommons.org/licenses/by/4.0/, as recorded in the arXiv OAI licence metadata for this exact version and shown on the abstract page https://arxiv.org/abs/2609.29748v1. Full rights notice: licenses/arxiv-2609.29748-CC-BY-4.0.txt.\par\smallskip
\noindent\textit{Editorial scope.} Counted unit: the original \texttt{theorem} environment \texttt{thm:main} at source lines 581--593 together with its complete proof at lines 595--861; the delivered document keeps source lines 142--862 so that every referenced label and every citation resolves inside it. Native editable source is the paper's own concatenated TeX; the AMS wrapper is editorial formatting only. No publisher class, upstream script or unrelated artwork is redistributed.
        \begin{equation}\label{axiom}
        N_G(x) = Ax^{\delta} + O(x^{\eta})\qquad(x\to\infty).
        \end{equation}
        It is proven in \cite[Chapter VI]{knopfmacher2015abstract} for any arithmetic semigroup $G$ satisfying Axiom A that
        \begin{equation}\label{primecounting}
    \pi_G(x)\sim\frac{x^\delta}{\delta\log x}\quad(x\rightarrow\infty).
        \end{equation}
        Quite a few classical examples of Prime Number Theorems are special cases of this result (e.g. the classical Prime Number Theorem and Landau's Prime Number Theorem \cite{landau2010einfuhrung} for ideals of bounded norm in rings of integers of number fields).

\subsection{The Dickman \texorpdfstring{$\rho$}{rho} function}
A lot is known about smoothness in the case where $G = \mathbb{N}_{\geq 1}$ under multiplication, $P$ is the set of primes and $|\cdot|$ is the usual absolute value.
The Dickman $\rho$ function describes the limiting proportion of smooth integers when $u=\log x/\log y$ is fixed. The same function appears in our asymptotic formula for $\Psi_G(x,y)$ under Axiom A. We record its definition and the properties needed in the proof of our main theorem.
\begin{definition}\label{def:dickman}
The Dickman function is the unique continuous function
$\rho:[0,\infty)\to\mathbb R$, differentiable on
$(1,\infty)$, satisfying
\[
\rho(u)=1\qquad(0\le u\le1)
\]
and
\[
u\rho'(u)+\rho(u-1)=0\qquad(u>1).
\]
\end{definition}

Integrating the defining differential equation gives,
for every integer $N\ge1$,
\begin{equation}\label{rho-recursion}
\rho(u)
=
\rho(N)-\int_N^u\frac{\rho(v-1)}{v}\,dv,
\qquad N\le u\le N+1.
\end{equation}
This identity determines $\rho$ successively on the
intervals $[N,N+1]$. In particular,
\[
\rho(u)=1-\log u\qquad(1\le u\le2).
\]

The function $\rho$ is strictly positive on $[0,\infty)$,
equals $1$ on $[0,1]$, and is strictly decreasing on
$(1,\infty)$. Hence, for every fixed $U>0$,
\[0<\rho(U)\le\rho(u)\le1
\qquad(0\le u\le U).\]
For the definition and proofs of these properties, see
Hildebrand-Tenenbaum
\cite[equations (1.5)-(1.6), p. 414,
and Lemma 2.5, p. 426]{hildebrand1993integers}.


\subsection{Auxiliary estimates}

We establish the estimates needed to prove our main theorem
uniformly for $0<u\le U$, for every fixed $U>0$.
The proof adapts the classical induction based on the
largest-prime-factor decomposition; see
\cite[Section 1, p. 416]{hildebrand1993integers}.


\begin{lemma}[Generalised Mertens' estimate]\label{lem:mertens}
There is a constant $B_G$ such that
\[
T(x):=\sum_{\substack{p\in P\\|p|\le x}}|p|^{-\delta}
=
\log\log x+B_G
+O\!\left(\frac1{\log x}\right)
\qquad(x\ge2).
\]
\end{lemma}

\begin{proof}
See Knopfmacher
\cite[Chapter VI, Lemma 2.5]{knopfmacher2015abstract}.
\end{proof}

\begin{lemma}[An integral estimate]\label{lem:integral}
For fixed $\alpha>0$ and $c>1$, as $x\to\infty$,
\[
\int_c^x\frac{t^{\alpha-1}}{\log t}\,dt
=
\frac{x^\alpha}{\alpha\log x}
+
O_{\alpha,c}\!\left(
\frac{x^\alpha}{(\log x)^2}
\right).
\]
\end{lemma}

\begin{proof}
Integration by parts gives
\[
\int_c^x\frac{t^{\alpha-1}}{\log t}\,dt
=
\frac{x^\alpha}{\alpha\log x}
-\frac{c^\alpha}{\alpha\log c}
+\frac1\alpha
\int_c^x\frac{t^{\alpha-1}}{(\log t)^2}\,dt.
\]
For $x\ge c^2$, we split the remaining integral at $\sqrt{x}$.
On $[c,\sqrt{x}]$,
\[
\int_c^{\sqrt{x}}
\frac{t^{\alpha-1}}{(\log t)^2}\,dt
\le
\frac1{(\log c)^2}
\int_c^{\sqrt{x}}t^{\alpha-1}\,dt
\ll_{\alpha,c}x^{\alpha/2}.
\]
On $[\sqrt{x},x]$,
\[
\int_{\sqrt{x}}^x
\frac{t^{\alpha-1}}{(\log t)^2}\,dt
\le
\frac4{(\log x)^2}
\int_{\sqrt{x}}^x t^{\alpha-1}\,dt
\ll_\alpha\frac{x^\alpha}{(\log x)^2}.
\]
Since
\[
x^{\alpha/2}
=o\!\left(\frac{x^\alpha}{(\log x)^2}\right),
\]
the result follows.
\end{proof}

\begin{lemma}[Weighted prime sums]\label{lem:weighted}
For the exponent $0\le\eta<\delta$ in Axiom A,
\[
\sum_{\substack{p\in P\\|p|\le x}}|p|^{-\eta}
\ll\frac{x^{\delta-\eta}}{\log x}
\qquad(x\ge2).
\]
\end{lemma}

\begin{proof}
Recall from Lemma \ref{lem:mertens} that
\[
T(t):=\sum_{\substack{p\in P\\|p|\le t}}|p|^{-\delta}
=\log\log t+B_G+O(1/\log t).
\]
Choose $q_0>1$ below every prime norm, which is
possible by local finiteness and since every prime has norm $|p| > 1$. Then $T(q_0)=0$.
By enlarging the implied constant on a fixed initial
interval, Mertens' estimate holds for all $t\ge q_0$.

Since $|p|^{-\eta}=|p|^{-\delta}|p|^{\delta-\eta}$,
partial summation with cumulative sum $T(t)$ and
weight $t^{\delta-\eta}$ gives, for $x\ge q_0$, then
\begin{equation}\label{abel}
\sum_{\substack{p\in P\\|p|\le x}}|p|^{-\eta}
=x^{\delta-\eta}T(x)-(\delta-\eta)\int_{q_0}^x
T(t)t^{\delta-\eta-1}\,dt.
\end{equation}
The lower-endpoint term vanishes because $T(q_0)=0$.

Substituting the Mertens' estimate into \eqref{abel}
and integrating the main terms by parts, we obtain
\begin{align*}
\sum_{\substack{p\in P\\|p|\le x}}|p|^{-\eta}
&=
q_0^{\delta-\eta}(\log\log q_0+B_G)
+\int_{q_0}^x\frac{t^{\delta-\eta-1}}{\log t}\,dt\\
&\quad+
O\!\left(
\frac{x^{\delta-\eta}}{\log x}
+\int_{q_0}^x\frac{t^{\delta-\eta-1}}{\log t}\,dt
\right).
\end{align*}
Here the terms $x^{\delta-\eta}\log\log x$ and
$B_Gx^{\delta-\eta}$ have cancelled exactly.
Since $\delta-\eta>0$, Lemma \ref{lem:integral} gives
\[
\sum_{\substack{p\in P\\|p|\le x}}|p|^{-\eta}
\ll\frac{x^{\delta-\eta}}{\log x}
\]
for sufficiently large $x$.
\end{proof}
In an arithmetic semigroup, distinct primes can have the same
norm. We use the inclusion-exclusion identity below to account for
this possibility.
For each $t$ that is the norm of a prime element, let
\[
d_t=\#\{p\in P:|p|=t\}.
\]
Note that all sums indexed by $t$ below run over the possible distinct norms
of prime elements.

\begin{lemma}[Largest-prime-norm decomposition]\label{lem:ie}
For $x,y\ge1$,
\begin{equation}\label{ie}
\Psi_G(x,y)
=
1+\sum_{t\le y}\sum_{k=1}^{d_t}
(-1)^{k-1}\binom{d_t}{k}\Psi_G(x/t^k,t).
\end{equation}
\end{lemma}

\begin{proof}
For each $t$ that is the norm of a prime element, define
\[
P_t=\{p\in P:|p|=t\},
\quad \text{and} \quad 
A_p(t)=\{a\in S_G(x,t):p\mid a\}
\quad \text{for}\,\, p\in P_t.
\]
These sets are finite and $\#P_t=d_t$. We have
\[
\bigcup_{p\in P_t}A_p(t)
=
\left\{
a\in G\setminus\{1\}:|a|\le x,\ 
\max_{\substack{p\in P\\p\mid a}}|p|=t
\right\}.
\]
Indeed, membership in $S_G(x,t)$ bounds every prime factor
norm by $t$, while divisibility by some $p\in P_t$ ensures
that the bound is reached.

Let $J\subseteq P_t$ be non-empty, where $\#J=k$.
Since the primes in $J$ are distinct, unique factorisation
gives
\[
\bigcap_{p\in J}A_p(t)
=
\left\{
a\in S_G(x,t):
\prod_{p\in J}p\mid a
\right\}.
\]
Consider the map
\[
S_G(x/t^k,t)\longrightarrow\bigcap_{p\in J}A_p(t),
\qquad
b\longmapsto b\prod_{p\in J}p.
\]
Since
\[
\left|\prod_{p\in J}p\right|=t^k,
\]
this map is well-defined. Conversely, every element $a$
of the intersection can be written uniquely as
\[
a=b\prod_{p\in J}p.
\]
The quotient satisfies
\[
|b|=\frac{|a|}{t^k}\le\frac{x}{t^k},
\]
and every prime divisor of $b$ has norm at most $t$.
Thus $b\in S_G(x/t^k,t)$, proving surjectivity and
uniqueness of the preimage. The map is therefore a
bijection. If $x/t^k<1$ (as, by convention, $\Psi_G(z, t) = 0$ for $z < 1$), both sets are empty.

Consequently,
\[
\#\bigcap_{p\in J}A_p(t)=\Psi_G(x/t^k,t).
\]
Finite inclusion-exclusion gives
\begin{align*}
\#\bigcup_{p\in P_t}A_p(t)
&=
\sum_{\varnothing\ne J\subseteq P_t}
(-1)^{\#J-1}
\#\bigcap_{p\in J}A_p(t)\\
&=
\sum_{k=1}^{d_t}
(-1)^{k-1}\binom{d_t}{k}\Psi_G(x/t^k,t),
\end{align*}
because $P_t$ has exactly $\binom{d_t}{k}$ subsets
of cardinality $k$.

Finally, every non-identity element has a uniquely
determined largest prime factor norm. Hence
\[
S_G(x,y)
=
\{1\}\sqcup
\bigsqcup_{t\le y}
\left(\bigcup_{p\in P_t}A_p(t)\right).
\]
By taking cardinalities and substituting the
inclusion-exclusion formula we prove \eqref{ie}.
\end{proof}

Note that when different prime elements have different norms, i.e.
$d_t=1$ for each $t$, then \eqref{ie} reduces to
\[
\Psi_G(x,y)
=
1+\sum_{\substack{p\in P\\|p|\le y}}
\Psi_G(x/|p|,|p|).
\]
For $G=\mathbb N_{\geq 1}$, this is the classical Buchstab identity
used in the classical smooth-number induction proof in \cite[Chapter 9, Theorem 9.3]{de2023analytic}. Formula \eqref{ie}
is its replacement when distinct primes can have equal norms.\\
To use \eqref{ie} in the induction, we must control
the correction terms with $k\ge2$, which arise when
distinct prime elements have the same norm.
The following estimate will bound their total
contribution over large prime norms.
\begin{lemma}[Equal-norm prime pairs]\label{lem:collision}
For $L\ge2$, we have that
\[
\sum_{t>L}d_t^2t^{-2\delta}
\ll\frac1{\log L}.\]
\end{lemma}

\begin{proof}
For $T\ge2$, formula \eqref{primecounting} gives
\[
\sum_{T<t\le2T}d_t
=
\pi_G(2T)-\pi_G(T) \leq \pi_G(2T) \ll \frac{{(2T)}^\delta}{\log 2T}
\ll\frac{T^\delta}{\log T}.
\]
Therefore
\begin{align*}
\sum_{T<t\le2T}d_t^2t^{-2\delta}
&\le
T^{-2\delta}\sum_{T<t\le2T}d_t^2\\
&\le
T^{-2\delta}
\left(\sum_{T<t\le2T}d_t\right)^2\\
&\ll\frac1{(\log T)^2}.
\end{align*}
Summing over the intervals $(2^jL,2^{j+1}L]$, we obtain
\[
\sum_{t>L}d_t^2t^{-2\delta}
\ll
\sum_{j=0}^{\infty}
\frac1{(\log L+j\log2)^2}.
\]
Since the summand is decreasing as a function of $j$,
the last sum is at most
\begin{align*}
\frac1{(\log L)^2}
+\int_0^\infty
\frac{ds}{(\log L+s\log2)^2}
&=
\frac1{(\log L)^2}
+\frac1{\log2\,\log L}\\
&\ll\frac1{\log L},
\end{align*}
since $\log L \geq \log 2$. This proves the result.
\end{proof}

\begin{lemma}[Stieltjes summation with the Dickman weight]
\label{lem:stieltjes}
For every fixed integer $N\ge2$, uniformly for
$N<\tilde u\le N+1$ as $x\to\infty$, we have that
\begin{align*}
&\sum_{\substack{p\in P\\
x^{1/\tilde u}<|p|\le x^{1/N}}}
|p|^{-\delta}
\rho\!\left(\frac{\log x}{\log|p|}-1\right)=
\int_N^{\tilde u}\frac{\rho(v-1)}v\,dv
+O\!\left(\frac1{\log x}\right).
\end{align*}
\end{lemma}

\begin{proof}
By the definition of $T$, the sum equals
\[
\int_{(x^{1/\tilde u},\,x^{1/N}]}
\rho\!\left(\frac{\log x}{\log t}-1\right)\,dT(t).
\]
The half-open interval agrees with the strict lower
endpoint in the prime sum.

By Lemma \ref{lem:mertens},
\[
\sup_{x^{1/\tilde u}\le t\le x^{1/N}}
|T(t)-\log\log t-B_G|
\ll\frac{\tilde u}{\log x}
\ll\frac1{\log x}.
\]
On this interval, the function
\[
t\longmapsto
\rho\!\left(\frac{\log x}{\log t}-1\right)
\]
is continuous and nondecreasing, taking values in $[0,1]$.
Also, $T(t)-\log\log t-B_G$ has bounded variation on this
compact interval.

Stieltjes integration by parts therefore gives
\begin{align*}
&\left|
\int_{(x^{1/\tilde u},\,x^{1/N}]}
\rho\!\left(\frac{\log x}{\log t}-1\right)
\,d\bigl(T(t)-\log\log t-B_G\bigr)
\right|\\
&\qquad\le
3\sup_{x^{1/\tilde u}\le t\le x^{1/N}}
|T(t)-\log\log t-B_G|\\
&\qquad\ll\frac1{\log x}.
\end{align*}
It follows that
\begin{align*}
&\sum_{\substack{p\in P\\
x^{1/\tilde u}<|p|\le x^{1/N}}}
|p|^{-\delta}
\rho\!\left(\frac{\log x}{\log|p|}-1\right)\\
&\qquad=
\int_{x^{1/\tilde u}}^{x^{1/N}}
\rho\!\left(\frac{\log x}{\log t}-1\right)
\frac{dt}{t\log t}
+O\!\left(\frac1{\log x}\right).
\end{align*}
With
\[
v=\frac{\log x}{\log t},
\qquad
\frac{dt}{t\log t}=-\frac{dv}{v},
\]
the endpoints $x^{1/\tilde u}$ and $x^{1/N}$ correspond
to $\tilde u$ and $N$, respectively. Hence
\[
\int_{x^{1/\tilde u}}^{x^{1/N}}
\rho\!\left(\frac{\log x}{\log t}-1\right)
\frac{dt}{t\log t}
=
\int_N^{\tilde u}\frac{\rho(v-1)}v\,dv.
\]
This proves the lemma.
\end{proof}


\section{The main theorem}\label{sec:main}
We now establish the Dickman asymptotic for smooth
elements of an arithmetic semigroup satisfying Axiom A,
uniformly when $u$ lies in a fixed bounded interval.
The proof follows the classical induction on the
smoothness parameter, using the largest-prime-norm
decomposition \eqref{ie}. The additional terms arising
from distinct primes with equal norms are controlled
by Lemma \ref{lem:collision}.



\begin{theorem}\label{thm:main}
Let $G$ be an arithmetic semigroup satisfying Axiom A
as in \eqref{axiom}. For every fixed $U>0$, as $x\to\infty$,
\begin{equation}\label{main}
\Psi_G(x,y)
=
Ax^\delta\rho(u)
+O\!\left(\frac{x^\delta}{\log x}\right),
\qquad
u=\frac{\log x}{\log y},
\end{equation}
uniformly for $y\ge2$ and $0<u\le U$.
\end{theorem}

\begin{proof}
We first prove \eqref{main} uniformly for $0<u\le2$.
We then proceed by induction, extending the uniform
estimate from $0<u\le N$ to $0<u\le N+1$.

\smallskip
\noindent\textbf{The range $0<u\le1$.}
Here $y=x^{1/u}\ge x$. Since every prime divisor of an
element of norm at most $x$ has norm at most $x$,
\[
\Psi_G(x,x^{1/u})=N_G(x)=Ax^\delta+O(x^\eta).
\]
Moreover,
\[
\frac{x^\eta}{x^\delta/\log x}
=
\frac{\log x}{x^{\delta-\eta}}
\longrightarrow0,
\]
because $\delta-\eta>0$. Since $\rho(u)=1$ for
$0<u\le1$, this proves \eqref{main} uniformly in this range.

\smallskip
\noindent\textbf{The range $1<u\le2$.}
Put $y=x^{1/u}$, so that $y\ge\sqrt{x}$.
An element of norm at most $x$ cannot be divisible by two prime
factors whose norms exceed $y$, as
their product would have norm greater than $y^2\ge x$.

Every element counted by $N_G(x)-\Psi_G(x,y)$ therefore
has a unique prime factor $p$ with $|p|>y$.
For a fixed such prime, division by $p$ provides a
bijection with the set of elements of norm at most $x/|p|$. Thus,
\[
\Psi_G(x,y)
=
N_G(x)
-\sum_{\substack{p\in P\\y<|p|\le x}}N_G(x/|p|).
\]

Using Axiom A for $x/|p|\ge1$ one obtains
\begin{align*}
\Psi_G(x,y)
&=
Ax^\delta
-Ax^\delta
\sum_{\substack{p\in P\\y<|p|\le x}}|p|^{-\delta}+
O\!\left(
x^\eta+
x^\eta
\sum_{\substack{p\in P\\y<|p|\le x}}|p|^{-\eta}
\right).
\end{align*}
By Lemma \ref{lem:weighted},
\[
x^\eta
\sum_{\substack{p\in P\\y<|p|\le x}}|p|^{-\eta}
\ll\frac{x^\delta}{\log x}.
\]
Note that the term $x^\eta$ is absorbed into the same bound.

Furthermore, Lemma \ref{lem:mertens} gives
\begin{align*}
\sum_{\substack{p\in P\\y<|p|\le x}}|p|^{-\delta}
&=
\log\log x-\log\log y
+O\!\left(\frac1{\log x}+\frac1{\log y}\right)\\
&=
\log u+O(1/\log x),
\end{align*}
uniformly for $1<u\le2$, since $\log y=(\log x)/u$.
Consequently,
\[
\Psi_G(x,x^{1/u})
=
Ax^\delta(1-\log u)
+O(x^\delta/\log x).
\]
As $\rho(u)=1-\log u$ on $[1,2]$, this establishes
\eqref{main} uniformly for $0<u\le2$.

\smallskip
\noindent\textbf{The induction step.}
Let $N\ge2$ be an integer. Suppose that
\begin{equation}\label{ih}
\Psi_G(x,x^{1/u})
=
Ax^\delta\rho(u)
+O\!\left(\frac{x^\delta}{\log x}\right)
\end{equation}
holds uniformly for $0<u\le N$, whenever
$x^{1/u}\ge2$. We prove the same estimate uniformly for
$N<\tilde u\le N+1$.

Subtracting \eqref{ie} at the thresholds $x^{1/N}$
and $x^{1/\tilde u}$ gives
\begin{equation}\label{subtract}
\Psi_G(x,x^{1/\tilde u})
=
\Psi_G(x,x^{1/N})-S_1-\sum_{i\ge2}S_i,
\end{equation}
where
\[
S_i
=
\sum_{x^{1/\tilde u}<t\le x^{1/N}}
(-1)^{i-1}\binom{d_t}{i}\Psi_G(x/t^i,t),
\]
and $\binom{d_t}{i}=0$ when $i>d_t$.
These sums are finite.

\smallskip
\noindent\emph{The terms with $i\ge2$.}
Axiom A implies $N_G(t)\ll t^\delta$ for $t\ge1$.
Together with our convention for arguments below $1$,
this yields
\[
\Psi_G(x/t^i,t)\ll x^\delta t^{-i\delta}.
\]
Hence,
\[\sum_{i\ge2}|S_i|
\ll
x^\delta
\sum_{x^{1/\tilde u}<t\le x^{1/N}}
\sum_{i=2}^{d_t}\binom{d_t}{i}t^{-i\delta}.
\]
By \eqref{primecounting},
\[
d_tt^{-\delta}
\le\pi_G(t)t^{-\delta}
\ll\frac1{\log t}.
\]
Since
\[
t>x^{1/\tilde u}\ge x^{1/(N+1)},
\]
we have that $d_tt^{-\delta}\le1$ uniformly over the
summation range for sufficiently large $x$.

Using the fact that $\binom{d_t}{i}\le d_t^i/i!$, we obtain
\begin{align*}
\sum_{i=2}^{d_t}\binom{d_t}{i}t^{-i\delta}
&\le
\sum_{i=2}^\infty\frac{(d_tt^{-\delta})^i}{i!}\\
&\le
d_t^2t^{-2\delta}\sum_{i=2}^\infty\frac1{i!}\\
&\ll d_t^2t^{-2\delta}.
\end{align*}
Applying Lemma \ref{lem:collision} with
$L=x^{1/\tilde u}$ therefore gives
\begin{equation}\label{error}
\sum_{i\ge2}|S_i|
\ll
\frac{x^\delta}{\log(x^{1/\tilde u})}
=
\frac{\tilde u x^\delta}{\log x}
\ll\frac{x^\delta}{\log x}.
\end{equation}

\smallskip
\noindent\emph{The principal sum $S_1$.}
Regrouping the terms according to norms of prime elements gives
\begin{equation}\label{regroup}
\begin{split}
S_1
&=
\sum_{x^{1/\tilde u}<t\le x^{1/N}}
d_t\Psi_G(x/t,t)\\
&=
\sum_{\substack{p\in P\\
x^{1/\tilde u}<|p|\le x^{1/N}}}
\Psi_G(x/|p|,|p|).
\end{split}
\end{equation}
For every prime in this range, we have that
\[
N-1
\le
\frac{\log(x/|p|)}{\log|p|}
=
\frac{\log x}{\log|p|}-1
<
\tilde u-1
\le N.
\]
Moreover,
\[
x/|p|\ge x^{1-1/N}\longrightarrow\infty,
\qquad
\log(x/|p|)\ge(1-1/N)\log x,
\]
uniformly over the range. Also,
$|p|>x^{1/(N+1)}\ge2$ for sufficiently large $x$.
Thus every summand lies in the range of the uniform
induction hypothesis.

Applying \eqref{ih} with norm bound $x/|p|$ and
smoothness threshold $|p|$, we obtain
\begin{align*}
\Psi_G(x/|p|,|p|)
&=
Ax^\delta|p|^{-\delta}
\rho\!\left(\frac{\log x}{\log|p|}-1\right)\\
&\quad+
O\!\left(
\frac{x^\delta|p|^{-\delta}}{\log x}
\right).
\end{align*}
By Lemma \ref{lem:mertens},
\begin{align*}
\sum_{\substack{p\in P\\
x^{1/\tilde u}<|p|\le x^{1/N}}}|p|^{-\delta}
&=
\log(\tilde u/N)+O(1/\log x)\\
&=O(1).
\end{align*}
The sum of the errors is therefore
$O(x^\delta/\log x)$, and
\begin{equation}\label{sone}
\begin{split}
S_1
&=
Ax^\delta
\sum_{\substack{p\in P\\
x^{1/\tilde u}<|p|\le x^{1/N}}}
|p|^{-\delta}
\rho\!\left(\frac{\log x}{\log|p|}-1\right)\\
&\quad+O(x^\delta/\log x).
\end{split}
\end{equation}

\smallskip
\noindent\emph{Completion of the induction.}
Lemma \ref{lem:stieltjes} applied to \eqref{sone} gives
\[
S_1
=
Ax^\delta
\int_N^{\tilde u}\frac{\rho(v-1)}v\,dv
+O(x^\delta/\log x).
\]
Combining this with \eqref{subtract}, \eqref{error},
and the induction hypothesis at $u=N$, we obtain
\begin{align*}
\Psi_G(x,x^{1/\tilde u})
&=
Ax^\delta
\left(
\rho(N)-\int_N^{\tilde u}\frac{\rho(v-1)}v\,dv
\right)\\
&\quad+O(x^\delta/\log x).
\end{align*}
By \eqref{rho-recursion}, the expression in parentheses
is $\rho(\tilde u)$. Hence
\[
\Psi_G(x,x^{1/\tilde u})
=
Ax^\delta\rho(\tilde u)
+O(x^\delta/\log x),
\]
uniformly for $N<\tilde u\le N+1$.
This completes the induction.

For any fixed $U>0$, finitely many induction steps cover the range
$0<u\le U$. Their constants can be absorbed into a single
constant depending on $G$ and $U$, proving \eqref{main}.
\end{proof}

\begin{thebibliography}{99}
\bibitem[de2023analytic]{de2023analytic} Jean-Marie De Koninck and Florian Luca, \emph{Analytic number theory: exploring the anatomy of integers}, vol. 134, American Mathematical Society, 2023.
\bibitem[hildebrand1993integers]{hildebrand1993integers} Adolf Hildebrand and G\'erald Tenenbaum, \emph{Integers without large prime factors}, J. Th\'eor.\ Nombres Bordeaux \textbf{5} (1993), no. 2, 411--484, DOI 10.5802/jtnb.101.
\bibitem[knopfmacher2015abstract]{knopfmacher2015abstract} John Knopfmacher, \emph{Abstract analytic number theory}, Courier Dover Publications, 2015.
\bibitem[landau2010einfuhrung]{landau2010einfuhrung} Edmund Landau, \emph{Einf\"uhrung in die elementare und analytische Theorie der algebraischen Zahlen und der Ideale} (1918), Kessinger Publishing, 2010 (German).
\end{thebibliography}
\end{document}
